\chapter{Personalized Expected Utility}
\label{ch-personalized-test}

This chapter
is based on 
Ref.\cite{ang-li-thesis}.

This chapter assumes
that the reader has already read 
Chapter \ref{ch-personalized} on 
Personalized Treatment Effects.
Whereas Chapter 
\ref{ch-personalized}
is concerned with
finding
bounds for
$PNS_z$,
this chapter
will find bounds for $EU_z$, which is
 called the
Personalized Expected Utility (PEU).


For $y_0, y_1\in \bool$, let
us denote  
the {\bf  conditional joint experimental (causal,
counterfactual) distribution} by:

\beq
P_{y_0, y_1|z}=
P(\rvy_0=y_0, \rvy_1=y_1|z)
\;.
\eeq
Suppose
we are given 
a {\bf Utility function}\footnote{Ref.\cite{ang-li-thesis}
refers to the utility function as the
benefit function.}

\beq
\alp_{y_0, y_1}: \bool^2\rarrow \RR
\eeq
Let\footnote{
The notation $\beta,\gamma, \theta, \delta$ 
won't be used  again in this chapter.
We  mention it here so the reader
 can translate to and from
our equations 
and the equations
 in Ref.\cite{ang-li-thesis}
where this
$\beta,\gamma, \theta, \delta$
notation is used.}

\beq
\begin{array}{lll}
\beta = \alp_{0,1} &(y_0=0, y_1=1)& \text{ compliers}
\\
\gamma = \alp_{1,1}&(y_0=1, y_1=1) & \text{ always takers} 
\\
\theta = \alp_{0,0}& (y_0=0, y_1=0)&\text{ never takers}
\\
\delta = \alp_{1,0}& (y_0=1, y_1=0)&\text{ defiers}
\end{array}
\eeq


\beq 
\xymatrix{
\ar@{-}[r]\ar@{-}[d]&x=0\ar@{-}[r]&x=1\ar[r]&
\\
y=0\ar@{-}[d]
&
\rvy_0=0\ar@{-}@/^1pc/[dr]|{\alp_{0,1}}
\ar@{-}[r]^{\alp_{0,0}}
&\rvy_1=0
\\
y=1\ar[d]
&
\rvy_0=1\ar@{-}@/_1pc/[ur]|{\alp_{1,0}}
\ar@{-}[r]_{\alp_{1,1}}
&
\rvy_1=1
\\
&
}
\eeq

Define the {\bf conditional
 (i.e., personalized) expected utility} by

\beq
EU_z = E_{\rvy_0,\rvy_1|z}[\alp_{\rvy_0,\rvy_1}]=\alp_{i,j}P_{i,j|z}
\eeq
where we are using the Einstein
 summation convention
(repeated indices are to be summed over)
for the indices $i,j\in \bool$.

Compare the definition of $EU_z$
with that of two other causal effects used 
in his book:

\beq
PNS_z = P_{0,1||z}=P(compliers|z)
\eeq

\beq
Uplift = ATE_z = E_{1|1,z}-E_{1|0,z}
\eeq
As we shall see,
$EU_z$ contains the information
in $PNS_z$ and $ATE_z$,  plus much more.

 
One can
find the stratum $z^*$
such that
$z^*=\argmax ATE_z$ (as is done A/B= RCT testing) or
 $z^*=\argmax_z EU_z$ or
$z^*=\argmax_z PNS_z$.
This is called the
{\bf unit selection problem}.
It's called ``unit selection"
rather than ``stratum selection" 
because once the  
stratum $z^*$ is found,
one can find a unit 
(i.e., individual) within
that stratum.







\section{Goal of PEU Theory}
Everything that we said 
in Chapter \ref{ch-personalized}
in the section entitled ``Goal 
of PTE Theory"
applies to PEU theory too,
if we just replace $PNS_z$
by $EU_z$.
As in Chapter \ref{ch-personalized}, 
we will consider two types of bounds (for
$EU_z$ instead of $PNS_z$): 
(1) Bounds for an unspecified bnet,
(2) Bounds for specific bnet families.


Here is
an explanation 
of why $EU_z$
varies within a bounded region,
something that might 
not be obvious to the beginner.
In  Fig.\ref{fig-peu-utility},
we represent the utility
function $\alp_{y_0, y_1}$
by a unit vector $\hat{\alp}$,
the probability distribution 
$P_{y_0, y_1|z}$ by a vector $\vec{P}$,
and $EU_z$ by the dot product 
$\vec{P}\cdot \hat{\alp}$.
When the probability
vector $\vec{P}$ varies within a bounded region
(shown in green),
this causes
the dot 
product $EU_z=\vec{P}\cdot \hat{\alpha}$
to vary within a bounded interval (also
shown in green).



\begin{figure}[h!]
\centering
\includegraphics[width=2.5in]
{personalized-test/utility.png}
\caption{Bounds (shown in green)
on the probability
vector $\vec{P}$
induce bounds (shown in green) on the dot 
product $EU_z=\vec{P}\cdot \hat{\alpha}$.
Here $\hat{\alp}$ is a unit vector
that stands for a normalized 
utility function. $\hat{\alp}$ 
does not vary but $\vec{P}$ does.
} 
\label{fig-peu-utility}
\end{figure}

\section{Bnets for PEU Theory}

Everything that we said 
in Chapter \ref{ch-personalized}
in the section entitled ``Bnets 
for PTE Theory"
applies to PEU theory too,
if we just replace $PNS_z$
by $EU_z$.




\section{Bounds on $EU$ for unspecified bnet}
Define the
{\bf balanced  utility} by
\beq
\alp_B = \alp_{0,0} + \alp_{1,1}
\;,
\eeq
the
{\bf unbalanced utility} by
\beq
\alp_U =
\alp_{1,0}+\alp_{0,1}
\;,
\eeq
and their difference by:

\beq
\s = \alp_U - \alp_B
\;.
\eeq
We will
also use the abbreviations:

\beq
\alp_{1-0, j}=\alp_{1,j} -\alp_{0,j}
\eeq
and

\beq
\alp_{j,1-0}=\alp_{j,1} -\alp_{j,0}
\;.
\eeq


\begin{claim}
In general, 

{\renewcommand\arraystretch{1.5}
\beq
\left\{
\begin{array}{ll}
\max p_{[1\ldots4]}
\leq
EU_z
\leq
\min p_{[5\ldots 8]}
&\text{ if } \s<0
\\
\max p_{[5\ldots 8]}
\leq
EU_z
\leq
\min p_{[1\ldots 4]}
&\text{ if } \s>0
\end{array}
\right.
\eeq
}
where\footnote{$p_{[a\ldots b]}=(p_a, p_{a+1}, \ldots p_b)$
for integers $a, b$ such that $a<b$.}

{\renewcommand\arraystretch{1.5}
\beq
\left\{
\begin{array}{c}
p_1=\alp_{0, 1-0}E_{1|1,z} +
\alp_{j,0}E_{j|0,z}
\\
p_2=
\alp_{0-1,1}E_{0|0,z}
+ \alp_{1,j}E_{j|1,z} 
\\
p_3 =p_5
+\s O_{*,*|z}
\\
p_4 = p_1
-\s E_{1|0,z}
+ \s (1-O_{*,*|z})
\end{array}
\right.
\eeq
}

{\renewcommand\arraystretch{1.5}
\beq
\left\{
\begin{array}{c}
p_5 =

\alp_{1, 1-0}E_{1|1,z} 
+\alp_{j,0}E_{j|0,z}
\\
p_6 = 
p_1
-\s E_{1|0,z}
\\
p_7 =p_5
-\s E_{1|0,z}
+\s P(\rvy=1|z)
\\
p_8 = p_1
-\s P(\rvy=1|z)
\end{array}
\right.
\eeq
}

\end{claim}
\proof
See Ref\cite{ang-li-thesis}
or use the MRP (Matrix Representation
of Probabilities) defined in Chapter \ref{ch-personalized}.
\qed

Later on, we will
see that under monotonicity,
these bounds collapse to a point 
estimate of $EU_z$.
But what if monotonicity doesn't hold?
In such cases, one can use
the midpoint
of the bounds
as a non-rigorous point estimate of $EU_z$.



\begin{claim}
In general,

\beq
P_{0,0|z}=E_{0|1,z}-P_{1,0|z}
\eeq

\beq
P_{1,1|z}=E_{1|0,z}-P_{1,0|z}
\eeq

\beqa
P_{0,1|z}
&=&
1-P_{0,0|z}-P_{1,1|z}-P_{1,0|z}
\\
&=&
 1-E_{0|1,z}-E_{1|0,z}+P_{1,0|z}
\eeqa


\end{claim}
\proof

\beq
P_{0,0|z}+P_{1,0|z}=
 P(\rvy_1=0|z)=E_{0|1,z}
\eeq

\beq
P_{1,1|z}+P_{1,0|z}=
 P(\rvy_0=1|z)=E_{1|0,z}
\eeq

\qed

Recall the 
definition of monotonicity. Monotonicity holds iff

\beq
P(\rvy_0=1, \rvy_1=0)=0
\;.
\eeq

\begin{claim}
\label{cl-peu-trivial-mono}
In general, 
\beq
EU_z = \alp_{0,0}P_{0,0|z}
+\alp_{1,1}P_{1,1|z}
+\alp_{0,1}P_{0,1|z}
+
\alp_{1,0}P_{1,0|z}
\eeq
Hence, if monotonicity holds,
or $\alp_{1,0}=0$, then

\beq
EU_z = \alp_{0,0}P_{0,0|z}
+\alp_{1,1}P_{1,1|z}
+\alp_{0,1}P_{0,1|z}
\eeq

\end{claim}
\proof Trivial
\qed

Claim \ref{cl-peu-trivial-mono}
is trivial, but
it provides a good motivation 
and inspiration for the
following less trivial claim.


\begin{claim}
\label{cl-peu-mono}
In general,

\beq
EU_z = \alp_{0,0} 
\underbrace{E_{0|1,z}}_{=1-E_{1|1,z}}
+ \alp_{1,1}E_{1|0,z}
+ \alp_{0,1}
\underbrace{( 1-E_{0|1,z}-E_{1|0,z})}
_{= E_{1|1,z}-E_{1|0,z}=ATE_z}
+ \s P_{1,0|z}
\eeq
Hence, if monotonicity holds,
or $\s=0$, then
\beq
(EU_z)_{\s=0} = \alp_{0,0} E_{0|1,z}
+ \alp_{1,1}E_{1|0,z}
+ \alp_{0,1}ATE_z
\eeq
\end{claim}
\proof
In general,
\beq
EU_z =
\left\{
\begin{array}{r}
\quad\alp_{0,0}P_{0,0|z}
\\
+\alp_{1,1}P_{1,1|z}
\\
+\alp_{0,1}P_{0,1|z}
\\
+\alp_{1,0}P_{1,0|z}
\end{array}
\right.
=
\left\{
\begin{array}{c}
\quad\alp_{0,0}(E_{0|1,z}-P_{1,0|z})
\\
+\alp_{1,1}(E_{1|0,z}-P_{1,0|z})
\\
+\alp_{0,1}(1-E_{0|1,z}-E_{1|0,z}+P_{1,0|z})
\\
+\alp_{1,0}P_{1,0|z}
\end{array}
\right.
\eeq
Hence,
\beq
EU_z = (EU_z)_{\s=0} + \s P_{1,0|z}
\eeq

\qed

Perhaps this will make Claim \ref{cl-peu-mono}
less mysterious to you.
Note that Claims \ref{cl-peu-trivial-mono}
and \ref{cl-peu-mono} imply the following: 

\beq
\left[\pder{EU_z}{\alp_{1,0}}\right]
_{\alp_{0,0}, \alp_{1,1}, \alp_{0,1}}=
\left[\pder{EU_z}{\s}\right]
_{\alp_{0,0}, \alp_{1,1}, \alp_{0,1}}=P_{1,0|z}
\eeq


\section{Bounds on $EU$
for specific bnet  families}


\begin{figure}[h!]
$$\xymatrix{
\rvz\ar[rrd]\ar[dd]
&\rva\ar[ddl]\ar[dr]
\\
&
&(\rvy_0, \rvy_1)\ar[d]
\\
\rvx
\ar[rr]
&&\rvy
}
$$
\caption{Bnet for $\rvx\rarrow \rvy$ with
 2 observed confounders $\rvz,\rva$. }
\label{fig-peu-2-confounders}
\end{figure}

Fig.\ref{fig-peu-2-confounders} gives
an example of a bnet that satisfies
$(\rvz\cup\rva)\cap de(\rvx)=\emptyset$.


\begin{claim}\label{cl-peu-2-confounders}
If $(\rvz\cup\rva)\cap de(\rvx)=\emptyset$, then

\beq
L\leq PNS_z \leq U \quad\text{(see Claim 
\ref{cl-minimal-bounds})}
\eeq

\beq
\left\{
\begin{array}{ll}
p_1 + \s U \leq EU_z \leq p_1+ \s L
& \text{ if } \s<0
\\
p_1 + \s L \leq EU_z \leq p_1+ \s U
& \text{ if } \s>0
\end{array}
\right.
\eeq
where

\beq
p_1=\alp_{0, 1-0}E_{1|1,z} +
\alp_{j,0}E_{j|0,z}
\eeq

\beq
L=\sum_{a}P(a|z)
\max\left\{
\begin{array}{c}
0
\\
E_{1|1,a,z}-E_{1|0,a,z}
\\
E_{0|0,a,z}-P(\rvy=0|a,z)
\\
E_{1|1,a,z}-P(\rvy=1|a,z)
\end{array}
\right \}
\eeq

\beq
U=
\sum_{a}P(a|z)
\min\left\{
\begin{array}{c}
E_{1|1,a,z}
\\
E_{0|0,a,z}
\\
O_{*, *|a,z}
\\
E_{*|*,a,z}-O_{*, *|a,z}
\end{array}
\right\}
\eeq
\end{claim}
\proof
\qed


In Chapter \ref{ch-personalized}, 
we defined what we mean when we say that
 $\rvz$ is a partial or pure mediator
relative to $(\rvx,\rvy)$.
Fig.\ref{fig-partial-mediator-2}
is a copy of 
Fig.\ref{fig-partial-mediator}.

\begin{figure}[h!]
$$\xymatrix{
*++[F-o]{\rvu}\ar[dd]\ar[r]
&(\rvz_0, \rvz_1)\ar[r]
&\rvz\ar[rd]
\\
&
&
&(\rvy_0, \rvy_1)\ar[dr]
\\
\rvx\ar[rruu]
\ar[rrrr]
&&&&\rvy
}$$
\caption{Example of a bnet in  
which $\rvz$ is a partial mediator
relative to $(\rvx, \rvy)$.
If $\rvz$ is a partial mediator
relative to $(\rvx, \rvy)$,
and, in addition, 
 there is no arrow
$\rvx\rarrow \rvy$,
we say $\rvz$
is a pure
mediator relative to $(\rvx,\rvy)$.
}
\label{fig-partial-mediator-2}
\end{figure}

\begin{claim} (Bounds if $\rvz$ 
is partial 
mediator)
\label{cl-peu-partial-med}

If $\rvz$ is a partial mediator
relative to $(\rvx, \rvy)$,
and $\rvx, \rvy, \rvz\in\bool$, then

\beq
L\leq PNS_z \leq U\quad \text{(see
 Claim \ref{cl-pte-partial-med} )}
\eeq

\beq
\left\{
\begin{array}{ll}
p_1 + \s U \leq EU_z \leq p_1+ \s L
& \text{ if } \s<0
\\
p_1 + \s L \leq EU_z \leq p_1+ \s U
& \text{ if } \s>0
\end{array}
\right.
\eeq
where

\beq
p_1=\alp_{0, 1-0}E_{1|1,z} +
\alp_{j,0}E_{j|0,z}
\eeq

\beq
L=
\max\left\{
\begin{array}{c}
0
\\
E_{1|1,z}-E_{1|0,z}
\\
E_{0|0,z}-P(\rvy=0|z)
\\
E_{1|1,z}-P(\rvy=1|z)
\end{array}
\right \}
\eeq

\beq
U= \min
\left\{
\begin{array}{c}
E_{1|1,z}
\\
E_{0|0,z}
\\
O_{*, *|z}
\\
E_{*|*,z}-O_{*, *|z}
\\
new
\end{array}
\right\}
\eeq


\beq
new=
\sum_{z_0, z_1}
\min
\left\{
\begin{array}{c}
O_{1|1,z_1,z}\\
 O_{0|0,z_0,z}
\end{array}
\right\}
\min
\left\{
\begin{array}{c}
P(z_1| z)
\\
P(z_0|z)
\end{array}
\right\}
\eeq
\end{claim}

\begin{claim} (Bounds if $\rvz$ 
is pure 
mediator)

If $\rvz$ is a pure mediator
relative to $(\rvx, \rvy)$, 
and $\rvx, \rvy, \rvz\in\bool$, then
Claim \ref{cl-peu-partial-med} for a
partial mediator is still
valid, except that
now ``new"
equals

\beq
new=
\sum_{z_0, z_1}
\indi(z_0\neq z_1)
\min
\left\{\begin{array}{c}
O_{1|\rvx=z_1,z}\\
O_{0|\rvx=z_0, z}
\end{array}\right\}
\min
\left\{\begin{array}{c}
P(\rvz=z_1| \rvx=1,z)
\\
P(\rvz=z_0| \rvx=0, z)
\end{array}\right\}
\eeq
\end{claim}